\chapter{States, Reversibility, and Kinematics}

Classical mechanics is not only a collection of equations for falling bodies,
orbiting planets, or charged particles.  Before choosing any particular
equation, we can ask a more general question: what must a law of motion look
like in order to count as a classical law?  We shall begin with systems so
simple that every possible history can be drawn on the blackboard.  Once the
ideas of state, determinism, reversibility, and conservation are clear there,
we shall return to particles moving in ordinary space.

\section{Two Kinds of Question}

There are two kinds of question in mechanics.  One asks for the specific law
that governs a particular system.  A planet moving in the gravitational field
of a massive body obeys one law; an electrically charged particle moving in a
magnetic field obeys another.  These are questions about the detailed laws of
nature.

The other question concerns the framework common to all such examples.  Are
there rules that every allowable law must satisfy?  Which features belong to
classical mechanics itself, independently of the force or the object under
study?  The specific examples will illustrate the answer, but it is this
second question that organizes the subject.

The most useful first example is an abstract coin on an abstract table.  We
retain only one property of the coin: it can show heads or tails.  Its space of
configurations is therefore
\begin{equation}
  \mathcal C=\{H,T\}.
\end{equation}
The word \emph{configuration} means one of the possible ways the system can
be.  An initial condition is a choice of one of them.  A one-sided coin would
be still simpler, but it would have only one configuration and nothing could
ever happen.  Two states are enough to make dynamics possible.

\section{The Coin: State, Law, and History}

Let time advance in integer steps, as if the world were illuminated by a
stroboscope:
\begin{equation}
  t=\ldots,-2,-1,0,1,2,\ldots .
\end{equation}
A law of motion is an update rule.  It assigns to the state at time \(t\) a
state at time \(t+1\).  The simplest rule does nothing:
\begin{equation}
  H\longmapsto H,
  \qquad
  T\longmapsto T.
  \label{eq:coin-identity}
\end{equation}
Starting from \(H\), the entire history is \(H,H,H,\ldots\); starting from
\(T\), it is \(T,T,T,\ldots\).  The rule is boring, but it is completely
predictive.

The other reversible rule exchanges the two states:
\begin{equation}
  H\longmapsto T,
  \qquad
  T\longmapsto H.
  \label{eq:coin-flip}
\end{equation}
The histories now alternate:
\begin{equation}
  H,T,H,T,\ldots,
  \qquad
  T,H,T,H,\ldots .
\end{equation}
An arrow points from the state before the update to the state after it.  Thus
the tail of an arrow says where the system starts and the arrowhead says where
it goes.

\begin{figure}[htbp]
  \centering
  \includegraphics[width=0.86\textwidth]{lecture_01_figure_02.png}
  \caption{The two configurations of the coin and the directed flip law on
  the lecture blackboard.  The clean update rule is given in
  Eq.~\eqref{eq:coin-flip}.}
  \lectureframe{00:06:50}
  \label{fig:coin-board}
\end{figure}

We can encode the faces by a two-valued variable,
\begin{equation}
  \sigma=
  \begin{cases}
    +1,&H,\\
    -1,&T.
  \end{cases}
\end{equation}
The letter \(\sigma\) is traditionally used for two-valued quantities in
physics, especially in descriptions of spin, although no knowledge of spin is
needed here.  In this notation the two rules become equations of motion:
\begin{align}
  \sigma(t+1)&=\sigma(t),
  &&\text{identity law},\label{eq:sigma-identity}\\
  \sigma(t+1)&=-\sigma(t),
  &&\text{flip law}.\label{eq:sigma-flip}
\end{align}
The variable does not add new physics.  It gives a compact label for the
configuration and turns a graph into an equation.

\begin{figure}[htbp]
  \centering
  \includegraphics[width=0.86\textwidth]{lecture_01_figure_03.png}
  \caption{The flip law written as
  \(\sigma(t+1)=-\sigma(t)\).  The boxed values \(\sigma=\pm1\) label heads and
  tails.}
  \lectureframe{00:08:58}
  \label{fig:sigma-board}
\end{figure}

For either rule, the initial state and the law determine the state at every
later time.  This statement presumes that the system is \emph{closed}.  If
someone reaches in and turns the coin, then the hand is an external influence
not included in our description.  We have not disproved determinism; we have
changed the system.  A classical law is intended to govern a closed system,
with all influences relevant to its evolution included.

There are four mathematical functions from a two-element set to itself.  The
two not written above send both inputs to \(H\), or both inputs to \(T\).
They are deterministic toward the future, but they erase the distinction
between two possible pasts.  When reversibility is imposed below,
Eqs.~\eqref{eq:coin-identity} and \eqref{eq:coin-flip} are exactly the two
allowed maps.

\section{A Die: Orbits and Conserved Quantities}

Replace the coin by a die.  For present purposes, the only important fact
about a die is that it has six configurations,
\begin{equation}
  \mathcal C=\{1,2,3,4,5,6\}.
\end{equation}
An initial condition is a choice of one number.  We could again let every
state remain fixed, but directed graphs are more useful for less trivial
rules.  Consider
\begin{equation}
  1\longmapsto2\longmapsto3\longmapsto4
  \longmapsto5\longmapsto6\longmapsto1.
  \label{eq:six-cycle}
\end{equation}
Starting at \(4\), for example, the history is

\[
  4,5,6,1,2,3,4,\ldots .
\]
The six states form one cycle.\footnote{Source note: during the first reading
of the arrows, one spoken phrase says that \(6\) returns to \(6\).  The
immediately following phrase, the drawn cycle, and the enumerated history all
give the intended update \(6\mapsto1\).}

\begin{figure}[htbp]
  \centering
  \includegraphics[width=0.86\textwidth]{lecture_01_figure_06.png}
  \caption{The single six-state cycle.  A state moves to the next state at
  every stroboscopic step and returns after six updates.}
  \lectureframe{00:12:40}
  \label{fig:six-cycle-board}
\end{figure}

A rule can look more complicated without being structurally different.  For
example,
\begin{equation}
  1\to2\to5\to3\to4\to6\to1
  \label{eq:relabeled-cycle}
\end{equation}
is still a single six-cycle.  If the vertices are placed around the page in
the order \(1,2,5,3,4,6\), it becomes the same picture as
Eq.~\eqref{eq:six-cycle}.  Labels such as \(1\) and \(5\) are names; the cycle
structure does not depend on them.

A genuinely different law divides the six states into two cycles, for
example
\begin{equation}
  1\to2\to3\to1,
  \qquad
  4\to6\to5\to4.
  \label{eq:two-cycles}
\end{equation}
There is still one die and one law, but a trajectory that begins on one cycle
can never reach the other.  We may attach a number \(Q\) to the cycles,
\begin{equation}
  Q=
  \begin{cases}
    0,&\text{on the cycle }\{1,2,3\},\\
    1,&\text{on the cycle }\{4,5,6\}.
  \end{cases}
\end{equation}
The state changes at every step while \(Q\) does not.

\begin{definition}
A conserved quantity is a function of the state whose value is unchanged by
the law of motion.
\end{definition}

The same idea allows one fixed point, one two-cycle, and one three-cycle.  The
conserved label says which orbit the system occupies; it need not identify the
state within that orbit.  Thus conservation is compatible with motion.  What
is preserved is not necessarily the complete configuration but some
distinction among possible configurations.

\begin{figure}[htbp]
  \centering
  \includegraphics[width=0.86\textwidth]{lecture_01_figure_07.png}
  \caption{Several disconnected cycles on the board: a fixed point, a
  two-cycle, and a three-cycle.  The cycle label is conserved even while the
  state moves around its cycle.}
  \lectureframe{00:18:40}
  \label{fig:conserved-cycles-board}
\end{figure}

For a finite reversible system, every trajectory is a cycle.  A single cycle
has no nonconstant conserved label of this elementary kind, because the
trajectory eventually visits every state.  Multiple cycles permit such a
label because evolution never carries a state from one component to another.

\section{What Must Be Included in a State?}

The coin examples used a one-step rule: the next face depended only on the
present face.  A question from the audience tests whether that restriction is
essential.

\begin{classroomqa}
\audiencequestion{Could one construct a sequence of heads and tails in which
the next result depends on more than the immediately visible face, perhaps on
one or several earlier results?\lecturetimestamp{00:19:27--00:21:23}}

\lecturerresponse{Yes.  In that case \(H\) or \(T\) alone is not the complete
state.  If one previous face is needed, enlarge the state space to
\(
  \{HH,HT,TH,TT\}.
\)
Each ordered pair records the present face and the preceding face.  The
evolution can then again be written as a one-step rule on this four-state
space.  If still more history is needed, include that information as well.}
\end{classroomqa}

This answer gives the operational meaning of state.

\begin{definition}
The state of a system at one time consists of all the information needed,
together with the law of motion, to predict its future.
\end{definition}

State is therefore not merely a photograph.  It is whatever complete set of
data makes the evolution definite.  Later, when we describe a Newtonian
particle, its position alone will not be enough; velocity will also belong to
the initial data.  The coin with memory is the simplest version of the same
principle.

\section{Infinite State Spaces}

An infinite number of states does not require infinitely many objects.  Take
one particle constrained to the integer sites of an infinite line.  Its state
is an integer
\[
  n\in\mathbb Z.
\]
One law moves it one site to the right at every step:
\begin{equation}
  n(t+1)=n(t)+1.
  \label{eq:lattice-one}
\end{equation}
This orbit never closes, but it passes through every integer if it is followed
through both future and past.  In that extended sense it plays the role of one
endless cycle.

Another law moves two sites at a time:
\begin{equation}
  n(t+1)=n(t)+2.
  \label{eq:lattice-two}
\end{equation}
Now even sites evolve only through even sites, and odd sites only through odd
sites.  Parity is conserved:
\begin{equation}
  Q_{\rm parity}(n)=n\bmod 2,
  \qquad
  Q_{\rm parity}(n(t+1))=Q_{\rm parity}(n(t)).
\end{equation}
The passage from finite to infinite configuration space permits histories
that never repeat, but it does not change the meanings of state, orbit, or
conservation.

\section{Determinism Is Not Yet Reversibility}

We now draw a rule that predicts the future uniquely but fails a second
classical requirement.  Give a three-sided coin the states heads, tails, and
edge, and define
\begin{equation}
  H\to T,
  \qquad
  T\to E,
  \qquad
  E\to T.
  \label{eq:nonreversible-map}
\end{equation}
Starting from \(H\), the history is
\[
  H,T,E,T,E,\ldots;
\]
starting from either \(T\) or \(E\), the subsequent alternation is equally
unambiguous.  There is exactly one outgoing arrow at each state, so the rule
is deterministic toward the future.

The difficulty appears when we try to reconstruct the past.  If the present
state is \(T\), the preceding state could have been \(H\) or \(E\).  Two
histories have merged.  Reverse every arrow and the defect becomes immediate:
from \(T\), the reversed graph offers two possible moves.  The reverse rule is
not deterministic.

\begin{figure}[htbp]
  \centering
  \includegraphics[width=0.86\textwidth]{lecture_01_figure_08.png}
  \caption{The nonreversible \(H,T,E\) law and representative histories.
  Both \(H\) and \(E\) feed into \(T\), so the present does not determine a
  unique past.}
  \lectureframe{00:28:50}
  \label{fig:nonreversible-board}
\end{figure}

\begin{definition}
A discrete law is reversible when the present state determines both a unique
next state and a unique previous state.  In a state graph, every state must
have exactly one outgoing arrow and exactly one incoming arrow.
\end{definition}

For a finite state space, a reversible update is a permutation of the states.
It never creates or destroys distinctions among states; it only rearranges
them.  This is why its graph decomposes into cycles.  The map in
Eq.~\eqref{eq:nonreversible-map} is a function, but not a one-to-one function,
and therefore has no inverse.

\begin{classroomqa}
\audiencequestion{Could a rule of this kind be allowed if the states
represented probabilities for a particle?\lecturetimestamp{00:29:46--00:31:17}}

\lecturerresponse{That would be a different framework.  Fundamental
classical mechanics, as formulated here, uses deterministic laws rather than
probabilistic update rules.  One may certainly invent and study such a rule,
but it is not a classical mechanical law in this sense.  Quantum mechanics is
not deterministic in the same way and nevertheless has an analog of
reversible evolution; that requires a later discussion.  Classical
statistical mechanics also uses probabilities, but it begins with underlying
mechanical laws and represents incomplete knowledge of their initial data.}
\end{classroomqa}

The distinction is important.  A probability distribution used because we
do not know the exact classical state is not an indeterministic law of motion.
It describes our information about a system whose microscopic evolution is
still taken to be deterministic and reversible.

\section{Exact Laws and Imperfect Knowledge}

Prediction requires two ingredients:
\begin{enumerate}
  \item the law of motion;
  \item the initial state.
\end{enumerate}
An exact law does not produce an exact forecast when the initial condition is
known only approximately.  For the coin, it is easy to imagine looking at the
table and learning the face perfectly.  Real mechanical variables are
continuous.  Positions and velocities can take arbitrary real values, and no
measurement supplies infinitely many decimal places.

Suppose the true initial state lies within a small region of possible states.
Under some dynamics that region remains small for a long time; under others,
nearby initial states separate rapidly.  In a chaotic system, small initial
uncertainties eventually grow into large differences.  The atmosphere is a
familiar example.  Better measurements can extend the useful forecast, but
for any fixed measurement accuracy there is a sufficiently distant time at
which the forecast becomes unreliable.

This practical limitation does not introduce ambiguity into the equations.
If the initial conditions were known exactly, a classical reversible law
would determine both future and past.  With finite precision, the quantitative
question is instead: how accurately must the initial state be known to make a
prediction for a specified length of time?

The graphical test now has a compact conclusion.  A law represented by one
incoming and one outgoing arrow at every state is deterministic and
reversible.  The corresponding statements for particles will be less
elementary, but they express the same classical structure.

\section{From Abstract States to Points in Space}

We next turn from coins and dies to point particles.  A particle may move in
one, three, five, or any other number of dimensions, but ordinary space gives
the immediate example.  To describe a point quantitatively, choose:
\begin{itemize}
  \item an origin;
  \item three mutually perpendicular axes;
  \item an orientation of those axes;
  \item a unit of length.
\end{itemize}
The origin is conventional.  It may be placed in Palo Alto or anywhere else,
and physical conclusions should not depend on that choice.  The axes may be
labeled \(x,y,z\), or \(x_1,x_2,x_3\).  After the \(x\)- and \(y\)-directions
have been fixed, the \(z\)-axis can point in either of two perpendicular
directions.  We choose the right-handed convention: with the thumb of the
right hand along \(x\) and the index finger along \(y\), the middle finger
points along \(z\).

Equal marks along the axes require a unit, such as meters, feet, inches, or
light-years.  Once these conventions are fixed, every point is labeled by
three numbers,
\begin{equation}
  (x,y,z).
\end{equation}
One can reach the point from the origin by displacements along the three axes.
A system of many particles requires one such triple for each particle.

\begin{figure}[htbp]
  \centering
  \includegraphics[width=0.86\textwidth]{lecture_01_figure_09.png}
  \caption{Cartesian axes and the three coordinates of a point.  The drawing
  fixes a coordinate convention; the point itself does not depend on where we
  choose to place the origin or axes.}
  \lectureframe{00:43:00}
  \label{fig:coordinates-board}
\end{figure}

\section{Vectors and Their Components}

A vector is an object with magnitude and direction.  The displacement from
the origin to a point is a vector, but vectors also describe velocity,
acceleration, electric fields, and many other quantities.  We shall write a
vector in boldface, such as \(\mathbf v\), and its magnitude as
\(\lvert\mathbf v\rvert\).  The magnitude is nonnegative and vanishes only for
the zero vector.

Geometrically, a vector is a free arrow.  Translating the arrow without
changing its length or direction does not change the vector.  In a particular
application a vector can be attached to a point, as when we specify the
electric field at a given location, but attachment is extra information and
not part of the abstract vector itself.

Relative to Cartesian axes, a vector has components
\begin{equation}
  \mathbf v=(v_x,v_y,v_z).
\end{equation}
The components are ordinary signed numbers, obtained by projection on the
axes.  They encode both magnitude and direction.  The three-dimensional
Pythagorean theorem gives
\begin{equation}
  \lvert\mathbf v\rvert^2
  =v_x^2+v_y^2+v_z^2,
  \qquad
  \lvert\mathbf v\rvert
  =\sqrt{v_x^2+v_y^2+v_z^2}.
  \label{eq:vector-magnitude}
\end{equation}
The ratios of the components determine the direction.  A negative component
means that the vector points partly along the negative direction of that
axis.

\subsection{Scalar multiplication and addition}

Multiplication by a real number changes a vector's length and, if the number
is negative, reverses its direction.  For \(c>0\), \(c\mathbf A\) points in
the direction of \(\mathbf A\) and has magnitude
\(c\lvert\mathbf A\rvert\).  In particular, \(-\mathbf A\) has the same
length as \(\mathbf A\) and points oppositely.

There are three equivalent pictures of vector addition.  First, place the
tail of \(\mathbf B\) at the head of \(\mathbf A\); the arrow from the tail of
\(\mathbf A\) to the head of \(\mathbf B\) is \(\mathbf A+\mathbf B\).
Second, place both tails together and complete a parallelogram; its diagonal
is the sum.  This form makes commutativity visible:
\begin{equation}
  \mathbf A+\mathbf B=\mathbf B+\mathbf A.
\end{equation}
Third, add corresponding components:
\begin{equation}
  C_i=A_i+B_i,
  \qquad i=x,y,z,
  \label{eq:component-addition}
\end{equation}
where \(\mathbf C=\mathbf A+\mathbf B\).  These rules also define arbitrary
linear combinations such as \(a\mathbf A+b\mathbf B\).

\begin{figure}[htbp]
  \centering
  \includegraphics[width=0.86\textwidth]{lecture_01_figure_10.png}
  \caption{The triangle and parallelogram constructions for vector addition.
  Both give the same vector \(\mathbf A+\mathbf B\).}
  \lectureframe{00:53:00}
  \label{fig:vector-addition-board}
\end{figure}

There is no general operation of dividing one vector by another.  There are,
however, two standard vector products.  The cross product will be left for a
later lecture.  Here we need the dot product.

\subsection{The dot product}

Project \(\mathbf A\) on the axis defined by \(\mathbf B\).  If the angle
between the vectors is \(\theta\), the signed component of \(\mathbf A\) along
\(\mathbf B\) is \(\lvert\mathbf A\rvert\cos\theta\).  Multiplying this
projection by \(\lvert\mathbf B\rvert\) defines
\begin{equation}
  \mathbf A\cdot\mathbf B
  =\lvert\mathbf A\rvert\,\lvert\mathbf B\rvert\cos\theta.
  \label{eq:dot-geometric}
\end{equation}
The result is a scalar, not another vector.  Equation~\eqref{eq:dot-geometric}
is symmetric in \(\mathbf A\) and \(\mathbf B\), so
\(\mathbf A\cdot\mathbf B=\mathbf B\cdot\mathbf A\).

\begin{figure}[htbp]
  \centering
  \includegraphics[width=0.86\textwidth]{lecture_01_figure_11.png}
  \caption{The blackboard construction of the dot product by projection.
  The typeset form is
  \(\mathbf A\cdot\mathbf B=AB\cos\theta\).}
  \lectureframe{00:58:30}
  \label{fig:dot-product-board}
\end{figure}

The sign carries geometric information.  The dot product is positive for an
acute angle, negative for an obtuse angle, and zero for perpendicular vectors.
Thus
\begin{equation}
  \mathbf A\perp\mathbf B
  \quad\Longleftrightarrow\quad
  \mathbf A\cdot\mathbf B=0
\end{equation}
for nonzero vectors.

In components,
\begin{equation}
  \mathbf A\cdot\mathbf B
  =A_xB_x+A_yB_y+A_zB_z.
  \label{eq:dot-components}
\end{equation}
This formula makes calculations immediate.  In particular,
\begin{equation}
  \mathbf A\cdot\mathbf A
  =A_x^2+A_y^2+A_z^2
  =\lvert\mathbf A\rvert^2.
  \label{eq:self-dot}
\end{equation}
Given the components of two vectors, Eqs.~\eqref{eq:dot-geometric} and
\eqref{eq:dot-components} determine the angle between them:
\begin{equation}
  \cos\theta
  =\frac{\mathbf A\cdot\mathbf B}
  {\lvert\mathbf A\rvert\,\lvert\mathbf B\rvert}.
\end{equation}

The law of cosines follows from the same algebra.  Let \(\mathbf A\) and
\(\mathbf B\) run from a common vertex to the other two vertices of a
triangle.  The third side is
\(\mathbf C=\mathbf A-\mathbf B\), not \(\mathbf A+\mathbf B\).  Therefore
\begin{align}
  C^2
  &=(\mathbf A-\mathbf B)\cdot(\mathbf A-\mathbf B)\notag\\
  &=A^2+B^2-2\mathbf A\cdot\mathbf B\notag\\
  &=A^2+B^2-2AB\cos\theta.
  \label{eq:law-of-cosines}
\end{align}

\begin{figure}[htbp]
  \centering
  \includegraphics[width=0.86\textwidth]{lecture_01_figure_04.png}
  \caption{The triangle used for the law of cosines.  The board correction
  identifies the closing side as
  \(\mathbf C=\mathbf A-\mathbf B\), leading to
  Eq.~\eqref{eq:law-of-cosines}.}
  \lectureframe{01:04:29}
  \label{fig:law-cosines-board}
\end{figure}

\section{Position, Velocity, and Acceleration}

Choose an origin and let \(\mathbf r\) be the position vector from the origin
to a particle.  Its components are the Cartesian coordinates:
\begin{equation}
  \mathbf r=(x,y,z),
  \qquad
  \lvert\mathbf r\rvert=\sqrt{x^2+y^2+z^2}.
\end{equation}
Changing the origin changes \(\mathbf r\), a question we shall postpone.  Once
the origin is fixed, motion is described by a time-dependent position,
\begin{equation}
  \mathbf r(t)=\bigl(x(t),y(t),z(t)\bigr).
\end{equation}
The trajectory is therefore encoded by three functions of time.

Velocity is the time derivative of position:
\begin{equation}
  \mathbf v(t)=\frac{d\mathbf r}{dt}
  =\dot{\mathbf r}(t),
  \qquad
  v_i=\frac{dx_i}{dt}=\dot x_i.
  \label{eq:velocity}
\end{equation}
A dot always denotes a derivative with respect to time.  The magnitude
\(\lvert\mathbf v\rvert\) is the speed.  Position and velocity are independent
directions in general: a particle at one location may be moving in any
direction allowed by the problem.

Acceleration is the rate of change of velocity and hence the second time
derivative of position:
\begin{equation}
  \mathbf a(t)=\frac{d\mathbf v}{dt}
  =\dot{\mathbf v}(t)
  =\frac{d^2\mathbf r}{dt^2}
  =\ddot{\mathbf r}(t),
  \qquad
  a_i=\ddot x_i.
  \label{eq:acceleration}
\end{equation}
Acceleration is also a vector.  If a particle moving in the positive
\(x\)-direction speeds up, \(a_x>0\).  If it slows down, \(a_x<0\).  What is
often called deceleration is simply acceleration opposite to the current
velocity.  More generally, velocity can change by changing its direction even
when its magnitude remains constant.

\begin{figure}[htbp]
  \centering
  \includegraphics[width=0.86\textwidth]{lecture_01_figure_05.png}
  \caption{Velocity components written as time derivatives of the position
  components.  In compact vector notation,
  \(\mathbf v=d\mathbf r/dt\).}
  \lectureframe{01:14:24}
  \label{fig:velocity-board}
\end{figure}

\subsection{Uniform acceleration on a line}

For one-dimensional motion, write
\begin{equation}
  x(t)=a+bt+ct^2.
  \label{eq:quadratic-motion}
\end{equation}
At \(t=0\), the particle is at \(x=a\).  Differentiating once and then twice
gives
\begin{align}
  v(t)=\dot x(t)&=b+2ct,\label{eq:line-velocity}\\
  a(t)=\ddot x(t)&=2c.\label{eq:line-acceleration}
\end{align}
Thus \(b\) is the initial velocity and \(2c\) is a constant acceleration.  If
\(c>0\), the velocity increases linearly; if \(c<0\), it decreases linearly.
This quadratic position law is the form taken by motion under a constant
acceleration, such as the idealized motion of a falling body near Earth's
surface.

\section{Uniform Circular Motion}

For the final example, let a particle move in the \(x\)-\(y\) plane on the
unit circle.  It begins on the positive \(x\)-axis at \(t=0\), and its angular
coordinate increases uniformly:
\begin{equation}
  \theta(t)=\omega t.
\end{equation}
The constant \(\omega\) is the angular frequency.  One revolution corresponds
to an increase \(2\pi\) in the angle, so the period is
\begin{equation}
  T=\frac{2\pi}{\omega}.
  \label{eq:circular-period}
\end{equation}
The angle may be allowed to increase beyond \(2\pi\); the position repeats
even though the continuously recorded angle does not reset.

On the unit circle,
\begin{equation}
  \mathbf r(t)
  =\bigl(\cos\omega t,\,\sin\omega t\bigr).
  \label{eq:circular-position}
\end{equation}
Differentiation gives
\begin{equation}
  \mathbf v(t)
  =\omega\bigl(-\sin\omega t,\,\cos\omega t\bigr).
  \label{eq:circular-velocity}
\end{equation}
The velocity is tangent to the circle.  We can prove its perpendicularity to
the radius without relying on the drawing:
\begin{align}
  \mathbf r\cdot\mathbf v
  &=\cos\omega t\,(-\omega\sin\omega t)
    +\sin\omega t\,(\omega\cos\omega t)\notag\\
  &=0.
  \label{eq:r-dot-v}
\end{align}

\begin{figure}[htbp]
  \centering
  \includegraphics[width=0.86\textwidth]{lecture_01_figure_12.png}
  \caption{Position and velocity components for motion on the unit circle.
  The cancellation in \(\mathbf r\cdot\mathbf v\) verifies that the velocity
  is tangent to the circle.}
  \lectureframe{01:24:30}
  \label{fig:circular-velocity-board}
\end{figure}

Differentiating once more,
\begin{equation}
  \mathbf a(t)
  =-\omega^2
  \bigl(\cos\omega t,\,\sin\omega t\bigr)
  =-\omega^2\mathbf r(t).
  \label{eq:centripetal-acceleration}
\end{equation}
The acceleration is parallel to the radius but points in the opposite
direction, toward the center.  This is the centripetal acceleration.\footnote{Terminology
note: the word ``centrifugal'' occurs at this point in the lecture, but the
equation and accompanying description identify the standard centripetal
acceleration, directed toward the center.}

\begin{figure}[H]
  \centering
  \includegraphics[width=0.86\textwidth]{lecture_01_figure_13.png}
  \caption{The component calculation of circular acceleration.  Each
  component is \(-\omega^2\) times the corresponding position component, so
  \(\mathbf a=-\omega^2\mathbf r\).}
  \lectureframe{01:26:30}
  \label{fig:circular-acceleration-board}
\end{figure}

Because the radius is one,
\begin{equation}
  \lvert\mathbf r\rvert=1,
  \qquad
  \lvert\mathbf v\rvert=\omega,
  \qquad
  \lvert\mathbf a\rvert=\omega^2.
\end{equation}
Increasing \(\omega\) increases the speed linearly and the required inward
acceleration quadratically.  For a circle of radius \(R\), the same calculation
would give \(v=\omega R\) and
\(\lvert\mathbf a\rvert=\omega^2R=v^2/R\); the lecture uses \(R=1\) to keep
the first calculation uncluttered.

We now have the elementary language needed for particle mechanics: states and
initial data, reversible evolution, coordinates, vectors, and the derivatives
that define velocity and acceleration.  The next task is to decide exactly
which initial data specify a particle and to find laws that carry those data
from one instant to the next.
